\documentclass[book.tex]{subfiles}

\begin{document}
% spherical harmonics -- r, theta, phi
% Each has mapping to specific basis polynomials
\chapter{Special Functions}
\label{chap:special_funcntions}
\noindent\rule{11cm}{0.4pt} \\
\textbf{Prerequisites:} \autoref{chap:calculus}, \autoref{chap:ode} \\
\textbf{Difficulty Level:} ** \\
\noindent\rule{11cm}{0.4pt}
\vspace{5mm}

Computations in physics are often facilitated by using special classes of functions that satisfy certain properties. While such functions will be intimately familiar to physicists, they are often more mathematically obscure to scientists from other disciplines. In this section, we will introduce some families of useful functions that we will have occasion to use later in the book.

\section{Spherical Harmonics}

Spherical harmonics form an orthogonal basis set of functions on the surface of the sphere. For this reason, spherical harmonics prove very useful in physics for decomposing radial functions to do Fourier analysis.

To define spherical harmonics, we start by first defining the associated Legendre polynomials. These polynomials are the solutions to the general Legendre equation

\begin{align}
(1 - x^2)\frac{d^2}{dx^2}P^m_{\ell}(x) - 2x\frac{d}{dx}P^m_{\ell}(x) + \left [\ell(\ell + 1) - \frac{m^2}{1-x^2} \right] P^m_{\ell}(x) = 0
\end{align}

Here the indices $m, \ell$ are both integers which are referred to respectively as the degree and order of the polynomial. In the special case the $m$ is $0$, the Legendre polynomials are given by Rodrigues' formula

\begin{align}
P_{\ell}(x) = \frac{1}{2^{\ell}\ell!}\frac{d^{\ell}}{dx^{\ell}} [(x^2 - 1)^{\ell}]
\end{align}


The spherical harmonics $Y^m_{\ell}(\theta, \varphi)$ are defined by
\begin{align}
Y^m_{\ell}(\theta, \varphi) &= N e^{im \varphi} P^m_{\ell}(\cos \theta)
\end{align}

Here $N$ is a normalization constant to satisfy the orthonormality constraint. 

\begin{figure}
    \centering
    % \includegraphics{figures/Mathematical Foundations/special_functions/spherical_harmonics.png}
    \includegraphics{figures/Mathematical Foundations/special_functions/spherical_harmonics1.png}
    \caption{A common visual representation of the first few real spherical harmonics. Blue portions represent regions where the function is positive, and yellow portions represent where it is negative. The distance of the surface from the origin indicates the absolute value of $Y_{\ell }^{m}(\theta ,\varphi)$ in angular direction $(\theta,\varphi)$
    \label{fig:spherical_harmonics} \url{https://en.wikipedia.org/wiki/Spherical_harmonics\#/media/File:Spherical_Harmonics.png}.}
\end{figure}

Alternatively, the spherical harmonics can be viewed as the eigenfunctions of the Laplace operator $\nabla^2$. They satisfy the equation

\begin{align}
r^2 Y^m_{\ell}(\theta, \varphi) &= -\ell (\ell + 1) Y^m_{\ell}(\theta, \varphi)
\end{align}

An arbitrary solution to Laplace's equation,
\begin{align}
\nabla^2 f &= 0
\end{align}
can then be written by the Fourier decomposition
\begin{align}
f(r, \theta, \varphi) &= \sum_{\ell = 0}^\infty \sum_{m = -\ell}^{\ell} f^m_\ell r^\ell Y^m_{\ell}(\theta, \varphi)
\end{align}

In fact, any square integrable function $f: S^2 \to \mathbb{C}$ on the sphere can be decomposed in this fashion. 

The product of two spherical harmonics can be decomposed as a sum of spherical harmonics using the Clebsch-Gordon coefficients $C^{LM}_{\ell_1 m_1 \ell_2 m_2}$
\begin{align}
Y^{m_1}_{\ell_1} Y^{m_2}_{\ell_2} &= \sum_{L, M} C^{L0}_{\ell_1 0 \ell_2 0} C^{LM}_{\ell_1 m_1 \ell_2 m_2} Y^M_{L}
\end{align}

The definition we have given above defines spherical harmonics as specific functions. There are a number of variants of the definition that can yield slightly different functions. For this reason, it is useful to talk of the space $H_\ell$ of all spherical harmonics of degree $\ell$.  This space proves useful since it allow for a group representation of $SO(3)$, the symmetry group of three dimensional rotation matrices. Suppose that $\rho$ is a rotation matrix and $\psi \in H_\ell$ is a spherical harmonic. Then
\begin{align}
\rho \cdot \psi &= \psi \circ \rho^{-1}
\end{align}
defines another spherical harmonic.

\section{Gamma Function}

The Gamma function $\Gamma$ provides a generalization of the factorial operation to the complex numbers. On the integers, the Gamma function has values

\begin{align}
\Gamma(n) &= (n-1)!
\end{align}

On the positive real numbers, the Gamma function can be defined by the integral
\begin{align}
\Gamma(r) &= \int_0^\infty t^{r-1}e^{-t} dt
\end{align}
This definition can be extended to the complex plane (except 0 and the negative integers) by analytic continuation. The plot below shows $\Gamma$ along the full real axis

\begin{figure}
    \centering
    \includegraphics{figures/Mathematical Foundations/special_functions/gamma_function.png}
    \caption{Plot of the values of the Gamma function along the real axis. Note that the Gamma function has poles at the negative integers.}
    \label{fig:my_label}
\end{figure}

\section{Bessel Functions}

The Bessel functions are the solutions to Bessel's differential equation
\begin{align}
x^2 \frac{dy^2}{dx^2} + x \frac{dy}{dx} + (x^2 - \alpha^2) y &= 0
\end{align}

They arise in the solution of Laplace's equation and Schrodinger's equation in cylindrical coordinates. Bessel functions can be thought of as a generalization of the sine functions. Since the Bessel functions arise as the solutions to a second order differential equation, there are two linearly independent familes of solutions which are typically called Bessel functions of the first type (denote $J_\alpha$) and Bessel functions of the second type (denoted $Y_\alpha$)

\begin{figure}
    \centering
    \includegraphics{figures/Mathematical Foundations/special_functions/bessel_first_kind.png}
    \caption{A plot of Bessel's functions of the first kind $J_\alpha$.}
    \label{fig:bessel_first}
\end{figure}

\begin{figure}
    \centering
    \includegraphics{figures/Mathematical Foundations/special_functions/bessel_second_kind.png}
    \caption{A plot of Bessel's functions of the second kind $Y_\alpha$.}
    \label{fig:bessel_second}
\end{figure}

\section{Airy Functions}

The Airy functions are given by the solutions to the differential equation
\begin{align}
\frac{d^2 y}{dx^2} - xy &= 0
\end{align}

The Airy functions are useful to describe the physics of rainbows and to solve Schrodinger's equation in triangular potential wells (which arise in semiconductor physics). Airy functions also come in two kinds, typically denoted $\mathrm{Ai}(x)$ and $\mathrm{Bi}(x)$. 

\begin{figure}
    \centering
    \includegraphics{figures/Mathematical Foundations/special_functions/airy_functions.png}
    \caption{Plots of $\mathrm{Ai}(x)$ and $\mathrm{Bi}(x)$}
    \label{fig:airy}
\end{figure}

\section{Exercises}
\begin{enumerate}
    \item Compute the explicit forms of the polynomials $P_1$, $P_2$, $P_3$.
    \item Compute the explicit forms of $Y^m_1$, $Y^m_2$, $Y^m_3$.
\end{enumerate}

\end{document}